% Source PDF page 33; manual page 2-4.
\section{Coordinate Systems}\label{sec:coordinate-systems}

When mathematically describing the motion of one object relative to others, it is
necessary to define one or more reference frames or coordinate systems. The
object's motion is written in terms of one of these systems. When dynamics are
involved, the relationship between those systems and an inertial, nonaccelerating
reference frame is also required.

Trajectory quantities such as position, velocity, and acceleration are stated in
terms of coordinate systems. A position near the earth's surface, for example, can
be described by Cartesian $x$, $y$, and $z$ components measured from the earth's
center, or by longitude, latitude, and altitude. The most useful representation
depends on the problem being solved, so TAOS provides several systems:

\begin{itemize}[leftmargin=1.6em]
  \item earth-centered and fixed, Cartesian (ECFC), used for the equations of
    motion and for all unit vectors;
  \item earth-centered, inertial, Cartesian (ECIC), used for inertial quantities;
  \item geocentric, used with a spherical earth for longitude and latitude;
  \item local geocentric horizon, used with a spherical earth for flight-path angles;
  \item geodetic, used with an ellipsoidal earth for longitude, latitude, and altitude;
  \item local geodetic horizon, used with an ellipsoidal earth for flight-path angles;
  \item body-fixed, used for Euler angles and vehicle angular orientation;
  \item velocity, used in defining bank angle;
  \item wind, used for aerodynamic angles;
  \item inertial-platform, used for inertial alignment platforms; and
  \item tangent-plane, used for range-safety and flat-earth calculations.
\end{itemize}

The equations of motion are written in ECFC coordinates, making ECFC the primary
TAOS coordinate system. Every coordinate-system unit vector is also expressed in
ECFC terms so that the vectors can be used directly for coordinate transformations.
The following subsections define the systems, their unit vectors, and their mutual
relationships. Functions that calculate unit vectors or perform coordinate transforms
are contained in the \texttt{coords} module.
